\chapter{The Beginnings of Planning}

The conception of a socialist planned economy to replace the market economy of capitalism was deeply embedded in Marxist thought, though little had been done by Marx or by his successors to develop it in detail. But the concept of planning was not specifically socialist; it was inherent in any reaction from the laissez-faire economy of the nineteenth century. The underlying theme of Witte's famous memorandum to the Tsar of 1899 was the need to plan the development of the Russian economy, though nothing was worked out precisely. The Bolsheviks in the crisis of the revolution and the civil war had no time for the theory of planning. But Lenin among others had been impressed by the degree to which the German war economy conformed to a pattern of centralized control and planning. Nor was this accidental. The final stage towards which capitalism was moving before the war by its own inner development was monopoly capitalism. By what Lenin called ``the dialectic of history'', war was hastening the transformation of monopoly capitalism into ``state monopoly capitalism'', which constituted ``the fullest material preparation for socialism''. ``Without the big banks'', wrote Lenin with emphasis in September 1917, ``socialism would be unrealizable''. The application of the German model to Russia presented all the difficulties inherent in the building of socialism in a backward economy. Though the recent growth of industry in Russia had been highly concentrated, and dependent directly or indirectly on the state, it was still in a primitive stage of organization, and had little theoretical or practical assistance or guidance to offer to socialist planners. But the principle of planning met with no resistance. The party programme of 1919 demanded ``one general state plan'' for the economy; and from this time party and Soviet resolutions on economic affairs regularly included the call for ``a single economic plan''. 
% 注：去除了段中 "banks'\" 和 "unrealizable''." 附近的杂乱引号。

For the present, however, plans for particular industries were a more promising approach. The most famous of these was the work of a commission for the electrification of Russia (Goelro) set up in February 1920. This plan had a special fascination for Lenin, who coined the aphorism: ``Communism is Soviet power, plus electrification of the whole country''. It was just a year later, on the eve of the introduction of NEP, that the decision was taken to establish a ``State General Planning Commission (Gosplan)''. But Lenin showed little enthusiasm at this time for current discussions of a general plan, which he dismissed as ``idle talk and pedantry''. While Goelro at once embarked on the practical task of planning and constructing a network of power stations, which later made an important contribution to the process of industrialization, Gosplan was confined for some years to academic exercises in comprehensive planning. Pronouncements were constantly made on the need for a single economic plan. But the party leaders, wedded to NEP and to priority for agriculture, remained lukewarm. From 1920 onwards the most active champions of planning were Trotsky and other critics of the official line. Planning was primarily a policy for industry, with remote and uncertain implications for agriculture; and its practical application meant more and more sweeping inroads into the market economy of NEP. In these conditions progress was slow. Plans for particular branches of production, including agriculture, were drawn up by the departments concerned. But these, unlike the Goelro plan, had no authority; nor was any attempt made to coordinate them. The president of Gosplan complained in the summer of 1924 that, three years after its foundation, there was still no ``single economic plan''. 
% 注：去除了 "imr", "Iwith", "[its" 等侧边乱码。

The case for comprehensive planning did not go unchallenged. Planning had been much discussed in general terms, but its practical implications had not been explored. A planned economy was a new and untried concept, which defied in hitherto unrecognized ways the traditional rules of a market economy. The aims of the planners were countered by formidable arguments drawn from the armoury of classical economics. Industry, and especially heavy industry, in the USSR was an inefficient high-cost producer; agriculture with its unlimited supply of peasant labour was a relatively low-cost producer. The maximum return on capital would be obtained by investing it in agriculture, by developing agricultural surpluses for export, and thus financing the import of industrial goods, including capital goods for the eventual development of industry. Even in the field of industrial production, in a country like the USSR where capital was scarce and unskilled labour superabundant, the rational course was to give priority to industries producing simple consumer goods which were labour-intensive, and not to industries producing capital goods which were capital-intensive. But a policy of priority for agriculture and for light consumer industry, however consonant with traditional economic analysis and with the principles of NEP, was the very antithesis of the ambition of the planners to speed the transformation of the USSR into a modern industrial country matching the industrial countries of the west. The arguments of the planners were political rather than economic, or perhaps belonged to a novel and unfamiliar kind of ``development economics''. The resistance to them, conscious and unconscious, of a large body of economists trained in the old school was strong and persistent. 
% 注：去除了 "(able" 中的单括号。

It was the scissors crisis of the autumn of 1923 which, by revealing the inadequacies of NEP, brought about measures of state intervention in the economy which were the first steps on the road to comprehensive planning. Wildly fluctuating prices disrupted orderly relations between country and town; heavy industry stagnated; figures of unemployment persistently rose. Price controls were introduced at the end of 1923. In January 1924 the party conference which called for a revival of the metal industry also, pursuing a perhaps unconscious train of thought, instructed Gosplan to ``establish a general perspective plan of the economic activity of the USSR for a number of years (five or ten)''. But the planners, though supported by Vesenkha as the champion of industry, still encountered the powerful opposition of Narkomzem and Narkomfin, the custodians of the market economy and of financial orthodoxy; and it was not till the following year that some advance was made. In August 1925 Gosplan issued its ``Control Figures of the National Economy'' (in essence, preliminary estimates) for the year beginning October 1, 1925. The figures were a mere outline, occupying with explanations and commentary less than a hundred pages; and they were marked by the resolute optimism which continued to inspire the efforts of Soviet planners. They were not mandatory; the economic departments were merely invited to take account of them in framing plans and programmes. The sceptics derided the figures as pure speculation. Sokolnikov, the People's Commissar for Finance, called the proposals for an increased emission of currency to finance the plan ``a formula of inflation''; and the excessive attention given to industry was attacked by Narkomzem. Of the party leaders, only Trotsky greeted the ``dry columns of figures'' with enthusiasm as ``the glorious music of the rise of socialism''. The other leaders received them with, at best, polite indifference. The difficulties of the grain collections after the 1925 harvest (see p. 79 above) thwarted the optimistic estimates of the planners and discredited their work. 
% 注：去除了 "uncons- cious" 的断字，修正了 "October i" 为 "October 1"，并清理了 "jthe", "jof" 标记。

In these circumstances it was not surprising that the important fourteenth party congress in December 1925, which ended in the defeat of Zinoviev and Kamenev, ignored the control figures, and had little to say about planning. Nevertheless, it was a turning-point. It was significant that Stalin should have attacked Sokolnikov as the principal advocate of keeping the USSR an agrarian country dependent on imports of industrial goods from abroad. The congress heralded the gradual abandonment by Stalin, once Zinoviev and Kamenev had been eliminated, of the peasant orientation inherent in NEP, and his conversion to far-reaching projects of industrialization. The congress resolution expressed determination to ``ensure the economic independence of the country, the development of the production of means of production, and the formation of reserves for economic manoeuvre''. All this, though its advocates may not have realized it, was a commitment to planning, and gave a strong stimulus to Gosplan and to the regional planning commissions which had been set up in many parts of the country. Hitherto the plans for particular industries and for agriculture had been drawn up by the departments concerned without any attempt at coordination. Now planning was to become comprehensive, for the economy as a whole. A new period had opened. 
% 注：去除了 ".dual" 的前导点。

The question was no longer whether to industrialize, but how to industrialize. Before 1925 industrial production had been climbing slowly out of the trough of the revolution and the civil war to reach its former levels. Hitherto the goal had been to restore what had been lost or destroyed. Progress in industrial technology in the capitalist countries since 1914 had actually enlarged the gap between the USSR and the industrialized countries of the west. New construction and new technological equipment were a crying need. Now that the way was clear for a fresh advance, major decisions of policy were required, which would have to be based on a broad plan for the whole economy. 

During the next two years the authority and prestige of Gosplan steadily increased. In March 1926, at a first Union planning congress, the task of Gosplan was divided into three branches---a ``general'' long-term plan, a ``perspective'' five-year plan, and annual operational plans; and a month later the party central committee, in a resolution on industrialization, demanded ``the reinforcement of the planning principle and the introduction of planning discipline''. The ``general plan'' proved to be an abortive enterprise. It was never completed, though it continued for some time to encourage visions of a long-term transformation of the economy. But the idea of planning for a five-year period caught the imagination, and stimulated the ambitions, of the planners. It compelled them to confine vague and remote prospects within the term of a fixed period; on the other hand, it was easier to produce optimistic estimates which were to be realized five years ahead than to limit oneself to the prospects of a single year. Alternative plans drafted by Gosplan and Vesenkha rivalled one another in their predictions of industrial development, and continued to excite controversy. The control figures of 1926 and 1927 were both fuller and more cautious than those of 1925. But interest in Gosplan shifted to the more ambitious projects of the five-year plan; and it came to be understood that the control figures were to be geared to this still hypothetical plan. The order was given to base the operational plans (called ``production and finance plans'') of particular industries on the estimates in the control figures. The structure of planning gradually took shape. 
% 注：去除了 "-r," 的边缘标记。

At this point a sharp division of opinion manifested itself in Gosplan between what were called the ``genetic'' and the ``teleological'' schools. It was significant that the former consisted mainly of non-party economists, most of them former Mensheviks, employed in Gosplan, and the latter of party members or of economists sensitive to the official party line. The ``geneticists'' argued that planning estimates must be based on the ``objective tendencies'' inherent in the economic situation, and were limited by those tendencies. The advocates of ``teleology'' maintained that the decisive factor in planning was the goal in view, and that one of its aims was to transform the economic situation and the tendencies inherent in it. Directives, not prediction, were the basis of the plan. This made planning a political, and not purely an economic, activity. Evidently both elements were present in all planning, and decisions rested on some balance or compromise between them. In practice the ``teleologists'' tended to reject the rules of a market economy, and claimed to override them by positive action; and this meant that they paid less attention to the conciliation of the peasant. These attitudes constituted, though this was rarely admitted, a direct challenge to NEP. In the later stages the effect of the ``teleological'' approach was to encourage the belief that, given sufficient determination and enthusiasm, no planning target was too high to attain. This mood came to dominate the preparation of the final version of the first five-year plan. 
% 注：去除了 "'members"、"—>"、"[them" 等边缘乱码，并修正了 "Stages" 为小写。

The identification of planning with industrialization was manifest from the outset. The underlying motive and driving-force was to develop Soviet industry, to catch up with the west, to make the USSR self-sufficient and capable of confronting the capitalist world on equal terms. An industry comparable to the industry of the western world had still to be created. The party congress in December 1925 accepted without question the principle of priority for the production of ``means of production'' over the production of consumer goods. This meant large investment in heavy industry which brought no immediate benefit to the consumer. To create reserves for investment from within industry itself, the costs of production were subjected to a ``regime of economy'', and brought within the scope of planning. Since, however, other opportunities of reducing costs were limited, the regime of economy bore most heavily on the workers; productivity must go up, or wages must come down. At the same time persistent attempts were made to force down retail prices by decree. But this resulted in increasing shortages of goods at the official prices, and left the consumer, especially in rural areas, at the mercy of the private trader, who still flourished under NEP conditions. The burden and discomforts of planning for industry began to come to the surface. 
% 注：去除了 "Ijnanifest", "torce" (修正为 force), "^est" (修正为 west), "Revolution]" 等排版错位，去除了 "same/time", "jirices" 中的乱码。

At first nobody tried to press these issues too hard. The costs of industrialization had not yet been fully counted. When Dzerzhinsky, the president of Vesenkha, died in July 1926, in the midst of a sharp controversy about the rate of investment in industry, his successor, Kuibyshev, revealed himself as a fervent advocate of what came to be known as ``forced industrialization''. Restraint was still imposed by the fact that the united opposition, Trotsky and Zinoviev, consistently pressed for more industrialization, and were at this time denounced by Stalin and Bukharin as ``super-industrializers''. What divided the two camps in the latter part of 1926 was not so much a difference about the desirability or the pace of industrialization, but the optimistic assumption of the majority, not shared by the opposition, that these policies could be pursued without severe strain on the economy, and in particular on its peasant sector. Opposition criticism was however silenced by the charge of lack of faith in the Soviet regime or in the working class. It was at this time that two major construction projects were approved---Dneprostroi, the great dam on the Dnieper river, and the Turksib railway connecting Central Asia with Siberia (see pp. 145--147 below). 
% 注：去除了 "JUU"、"difTerence" 的错误拼写。

The optimism of the last months of 1926 was succeeded by the anxieties of the following spring and summer, when the hostility of the west seemed to threaten the USSR with blockade or war. But this reversal of mood, far from calling a halt to the haste for industrialization, strengthened the determination to make the USSR self-sufficient, and to enable it to face a hostile capitalist world. Successive draft plans were prepared and circulated; and those who protested that the targets set were unrealistically high were soon outnumbered by those who called for more rapid and intensive progress. The ``regime of economy'' was followed by a campaign for the ``rationalization of production''---a term which covered a variety of pressures on workers and managers to increase efficiency and lower costs. ``Rationalization'' in several different forms could raise the productivity of labour, i.e. the output per worker employed. It could do so by tightening up organization, in management or on the shop-floor, by concentrating production in the most efficient units, by standardizing production and reducing the number of models produced. It could do so by arranging for the more efficient and economical utilization of existing plant and machinery. It could do so, above all, by modernizing and mechanizing the processes of production, in which the USSR lagged very far behind the major industrial countries. All these methods of rationalization were extensively tried from 1926 onwards, and achieved some success in reducing costs. But in a country like the USSR with scarce capital resources their scope was limited. In particular, the mechanization of industry, the largest source of rationalization, depended in this period mainly on the importation of machines from abroad, and very often on the employment of foreign personnel to give instruction in manning them. These conditions meant that the productivity of labour depended to a greater extent in the USSR than in the west on the physical energy of the workers. Productivity had to be raised primarily by harder, more efficient, better disciplined physical work; and every form of persuasion and pressure was directed to ensure this result. 
% 注：去除了 "—7", "fT TTfF trTTfn^", "fr^'^^^", "'^", "Xu^*/*-!", "t*^^J" 等等极端密集的边缘涂鸦乱码。

The implications of planning for other sectors of the economy were also disquieting, and were faced with reluctance. The cult of the peasant, wholeheartedly upheld by Bukharin, was still powerful throughout 1927; and the influence of Narkomzem, though now on the wane, still put a brake on the aspirations of the planners. Narkomfin continued to resist the assumption that unlimited credits could be made available for industrial expansion, and waged a stubborn struggle for what was dubbed ``the dictatorship of finance'' against ``the dictatorship of industry''. This raised the issue of ``financial'' controls, through manipulation of the supply of credit and the monetary system, as against ``physical'' controls through state decrees, which were exemplified in the heavy industries working directly on state orders, and in the monopoly of foreign trade. It was only gradually, even in the minds of the planners, that the financial instruments were recognized as inadequate, and replaced by direct ``physical'' controls. These controversies turned ultimately on the attitude to be adopted to the market economy which was the foundation of NEP. It was at first assumed that the planners would work within the market economy. Slowly and painfully, the incompatibilities between planning and industrialization, on the one hand, and NEP and the market economy on the other, came to light.
% 注：此处修复了严重的多栏断句倒置 ("[available for industrial expansion, and waged a stubborn strug- Narkomzem... gle")，现已拼合为 "Narkomfin continued to resist the assumption that unlimited credits could be made available for industrial expansion, and waged a stubborn struggle..."